\section{Reflection as Turning Over the Receiver}
\label{app:yaw_rotation}

Reflection can also be modeled as turning over the receiver, which is equivalent to rotating the reference frame depicted in \Fig{polarization_ellipse}
by $180\deg$ about a transverse axis in the $x$-$y$ plane.
This negates the $z$ axis and results in a two-dimensional reflection in the $x$-$y$ plane 
through the rotation axis.
%
Let $\hat{\Pv{a}}=(\cos\alpha,\sin\alpha,0)^T$ define the three-dimensional 
rotation axis in the physical space of \Fig{polarization_ellipse}\  
and use Rodrigues' Rotation Formula to express the $3\times3$ matrix 
for a 180$\deg$ rotation about $\hat{\Pv{a}}$.
\begin{equation}
\left( \begin{array}{ccc}
2a_x^2 - 1 & 2a_x a_y    & 0 \\
2a_x a_y   & 2a_x^2 - 1  & 0 \\
0          & 0           & -1 \\
  \end{array} \right)
  \end{equation}
In the $x$-$y$ plane, this transformation reduces to 
\begin{equation}
% W for weerspiegelen
\JM{W} = \left( \begin{array}{cc}
\cos2\alpha  & \sin2\alpha  \\
\sin2\alpha  & -\cos2\alpha 
  \end{array} \right)
  = \cos2\alpha\, \pauli{1} + \sin2\alpha\, \pauli{2},
\end{equation}
%
which is a linear combination of the two-dimensional reflection basis matrices 
defined by the two real-valued Pauli matrices.\footnote{
%
Both \pauli{1}\ and \pauli{2}\ 
satisfy the criteria of a $2\times2$ reflection matrix, which must be orthogonal 
($\JM{W}\JM{W}^T = \JM{W}^T\JM{W} = \pauli{0}$) and
have eigenvalues, $\lambda = \pm 1$, and determinant, $|\JM{W}|=-1$.
%
}
A reflection in the $x$-$y$ plane negates the \emph{handedness}\footnote{
%
Strictly, handedness is a property of only three-dimensional
vector spaces; this property is formally known as the orientation of the ordered basis. }
of the two-dimensional basis, which can also be reversed by swapping the cables 
used to propagate the orthogonally polarized signals.
%
In \citet{vmjr10}, the \emph{feed hand} defines the handedness of the receptor basis
and the reflection is performed through the axis defined by the \emph{symmetry angle}.

The polar decomposition of a two-dimensional reflection through the
axis defined by $(\cos\alpha, \sin\alpha)$ yields a unitary matrix parameterized by 
$\hat{\Pv{n}} = (\cos2\alpha\,\sin2\alpha,0)^T$,
%
\begin{equation}
\JM{W}(\alpha) = -\Ci\, \vRotation(\pm\pi/2).
\end{equation}

\begin{proof}
Note that $|\JM{W}|^{1/2}=\pm\Ci$ and
\begin{align*}
\JM{W}(\alpha)
    &= -\Ci \, \vRotation(\pi/2) \\
    &= -\Ci \left( \pauli{0}\cos\frac{\pi}{2} 
        + \Ci \, \Pv{n} \cdot \Pv{\sigma} \sin\frac{\pi}{2} \right) \\
    &= \Pv{n} \cdot \Pv{\sigma} \\
    &= \cos2\alpha\, \pauli{1} + \sin2\alpha\, \pauli{2}
\end{align*}
\end{proof}
%
\noindent
Congruence transformation by $\JM{W}(\alpha)$ rotates $\Pv{S}$ by $180\deg$ about the $\hat{\Pv{n}}$ axis defined by $\alpha$, which negates both the ellipticity angle $\chi$ and the position angle $\psi$.

In a basis defined by linearly-polarized receptors, the nominal axis of symmetry
has a position angle of $45\deg$.
The reflection defined by $\JM{W}(\pi/4)=\pauli{2}$ swaps $e_x$ and $e_y$.
%
With reference to \eqns{StokesQ}{through}{StokesV}, swapping $e_x$ and $e_y$
negates Stokes Q and V.
%
This negation is equivalent to a $\pm 180\deg$ rotation of the Stokes polarization vector about $\hat{\Pv{n}} = (0,1,0)^T$,
which is also the result of congruence transformation by $\JM{W}(\pi/4)$.

In the circular basis, the nominal axis of symmetry
has a position angle of $0\deg$, and $\JM{W}(0)=\pauli{1}$ negates $e_y$.
%
Negating either component 
of the electric field vector reverses the signs of both Stokes U and V, 
which is equivalent to the $\pm 180\deg$ rotation of the Stokes polarization vector about 
$\hat{\Pv{n}} = (1,0,0)^T$ under congruence transformation by $\JM{W}(0)$.
